\documentclass[10pt,a4paper]{amsart}
\usepackage{amsmath,amssymb,amsthm,graphicx}
% Retained minimal preamble for the Stacks Project chapter more-algebra.tex (More on Algebra).
% The chapter file pinned-original-source.tex carries only \input{preamble} on its first line, so no
% preamble text lives inside the chapter itself.  The Stacks Project uses one single shared preamble for
% every chapter; the snapshot stages it identically as sources/stacks-chow/preamble.tex and
% sources/stacks-etale-cohomology/preamble.tex.  The block below is copied mechanically, line for line,
% from sources/stacks-chow/preamble.tex lines 14-18 (the Xy-pic setup) and 154-217 (the native theorem
% declarations together with the mathematical macros).  Nothing was added, removed or redefined; the
% wrapper class and the amsmath/amssymb/amsthm/graphicx base packages are supplied by the assembler.
% Chapter: more-algebra.tex -- More on Algebra.
% For commutative diagrams we use Xy-pic
\usepackage[all]{xy}

% We use 2cell for 2-commutative diagrams.
\xyoption{2cell}

% Theorem environments.
%
\theoremstyle{plain}
\newtheorem{theorem}[subsection]{Theorem}
\newtheorem{proposition}[subsection]{Proposition}
\newtheorem{lemma}[subsection]{Lemma}

\theoremstyle{definition}
\newtheorem{definition}[subsection]{Definition}
\newtheorem{example}[subsection]{Example}
\newtheorem{exercise}[subsection]{Exercise}
\newtheorem{situation}[subsection]{Situation}

\theoremstyle{remark}
\newtheorem{remark}[subsection]{Remark}
\newtheorem{remarks}[subsection]{Remarks}

\numberwithin{equation}{subsection}

% Macros
%
\def\lim{\mathop{\mathrm{lim}}\nolimits}
\def\colim{\mathop{\mathrm{colim}}\nolimits}
\def\Spec{\mathop{\mathrm{Spec}}}
\def\Hom{\mathop{\mathrm{Hom}}\nolimits}
\def\Ext{\mathop{\mathrm{Ext}}\nolimits}
\def\SheafHom{\mathop{\mathcal{H}\!\mathit{om}}\nolimits}
\def\SheafExt{\mathop{\mathcal{E}\!\mathit{xt}}\nolimits}
\def\Sch{\mathit{Sch}}
\def\Mor{\mathop{\mathrm{Mor}}\nolimits}
\def\Ob{\mathop{\mathrm{Ob}}\nolimits}
\def\Sh{\mathop{\mathit{Sh}}\nolimits}
\def\NL{\mathop{N\!L}\nolimits}
\def\CH{\mathop{\mathrm{CH}}\nolimits}
\def\proetale{{pro\text{-}\acute{e}tale}}
\def\etale{{\acute{e}tale}}
\def\QCoh{\mathit{QCoh}}
\def\Ker{\mathop{\mathrm{Ker}}}
\def\Im{\mathop{\mathrm{Im}}}
\def\Coker{\mathop{\mathrm{Coker}}}
\def\Coim{\mathop{\mathrm{Coim}}}

% Boxtimes
%
\DeclareMathSymbol{\boxtimes}{\mathbin}{AMSa}{"02}

%
% Macros for moduli stacks/spaces
%
\def\QCohstack{\mathcal{QC}\!\mathit{oh}}
\def\Cohstack{\mathcal{C}\!\mathit{oh}}
\def\Spacesstack{\mathcal{S}\!\mathit{paces}}
\def\Quotfunctor{\mathrm{Quot}}
\def\Hilbfunctor{\mathrm{Hilb}}
\def\Curvesstack{\mathcal{C}\!\mathit{urves}}
\def\Polarizedstack{\mathcal{P}\!\mathit{olarized}}
\def\Complexesstack{\mathcal{C}\!\mathit{omplexes}}
% \Pic is the operator that assigns to X its picard group, usage \Pic(X)
% \Picardstack_{X/B} denotes the Picard stack of X over B
% \Picardfunctor_{X/B} denotes the Picard functor of X over B
\def\Pic{\mathop{\mathrm{Pic}}\nolimits}
\def\Picardstack{\mathcal{P}\!\mathit{ic}}
\def\Picardfunctor{\mathrm{Pic}}
\def\Deformationcategory{\mathcal{D}\!\mathit{ef}}

% --- appended 2026-10-12 (Chief): the project own definitions of the reference/slogan/history
% display environments, copied mechanically from the shared Stacks preamble lines 8-12
% (\usepackage{verbatim} plus the three \newenvironment lines, which the chapter preamble keeps
% verbatim from the project preamble). Nothing added, removed or redefined: these environments
% carry no mathematical content and the project itself typesets them as comments. Without these
% three lines a counted statement whose own text uses them fails with Environment slogan undefined.
\usepackage{verbatim}
\newenvironment{reference}{\comment}{\endcomment}
\newenvironment{slogan}{\comment}{\endcomment}
\newenvironment{history}{\comment}{\endcomment}
\usepackage{cleveref}
\crefname{originallemma}{original lemma}{original lemmas}
\Crefname{originallemma}{Original lemma}{Original lemmas}
\makeatletter
\newlabel{algebra-lemma-colimit-category-fp-modules}{{05N7}{source}{}{sourcepointer-1}{}}
\newlabel{algebra-lemma-colimit-category-fp-modules@cref}{{[originallemma][0][]05N7}{[1][0][]source}{}{}{}}
\makeatother
\crefname{originallemma}{original lemma}{original lemmas}
\Crefname{originallemma}{Original lemma}{Original lemmas}
\makeatletter
\newlabel{algebra-lemma-directed-colimit-exact}{{00DB}{source}{}{sourcepointer-2}{}}
\newlabel{algebra-lemma-directed-colimit-exact@cref}{{[originallemma][0][]00DB}{[1][0][]source}{}{}{}}
\makeatother
\crefname{originallemma}{original lemma}{original lemmas}
\Crefname{originallemma}{Original lemma}{Original lemmas}
\makeatletter
\newlabel{algebra-lemma-finite-projective}{{00NX}{source}{}{sourcepointer-3}{}}
\newlabel{algebra-lemma-finite-projective@cref}{{[originallemma][0][]00NX}{[1][0][]source}{}{}{}}
\makeatother
\crefname{originallemma}{original lemma}{original lemmas}
\Crefname{originallemma}{Original lemma}{Original lemmas}
\makeatletter
\newlabel{algebra-lemma-module-map-property-in-colimit}{{05LI}{source}{}{sourcepointer-4}{}}
\newlabel{algebra-lemma-module-map-property-in-colimit@cref}{{[originallemma][0][]05LI}{[1][0][]source}{}{}{}}
\makeatother
\crefname{originallemma}{original lemma}{original lemmas}
\Crefname{originallemma}{Original lemma}{Original lemmas}
\makeatletter
\newlabel{derived-lemma-morphisms-from-projective-complex}{{064B}{source}{}{sourcepointer-5}{}}
\newlabel{derived-lemma-morphisms-from-projective-complex@cref}{{[originallemma][0][]064B}{[1][0][]source}{}{}{}}
\makeatother
\crefname{originallemma}{original lemma}{original lemmas}
\Crefname{originallemma}{Original lemma}{Original lemmas}
\makeatletter
\newlabel{lemma-K-flat-resolution}{{14380}{source}{}{sourcepointer-6}{}}
\newlabel{lemma-K-flat-resolution@cref}{{[originallemma][0][]14380}{[1][0][]source}{}{}{}}
\makeatother
\title{The triangulated category of perfect complexes over a filtered colimit of rings is the colimit of the categories over the stages}
\author{The Stacks Project Authors}
\date{stacks-project:more-algebra.tex:local-snapshot-20261012; GNU FDL 1.2 or later}
\begin{document}
\maketitle
\setcounter{subsection}{1}
\noindent\textit{Editorial note.} All delivered mathematical text is the author's own, copied from the pinned Stacks Project chapter file more-algebra.tex; no statement or proof step was rewritten, reordered or removed.\par\medskip
\noindent\textit{Editorial note.} The chapter file pinned-original-source.tex carries only \textbackslash{}input\{preamble\} on its first line, so no preamble text lives inside the chapter. The retained file sources/stacks-more-algebra/retained-preamble.tex is a mechanical, line-for-line copy of the project's single shared preamble as staged in the snapshot; nothing was added, removed or redefined and no package was added.\par\medskip
\noindent\textit{Editorial note.} The source URL is the source-level Stacks Project URL, as required when more than one card is produced from one source; the per-claim tag is carried by claim\_key.\par\medskip
\noindent\textit{Editorial note.} The six references inside the counted proof span are recorded as non-load-bearing navigation pointers.\par\medskip
\noindent\textit{Allowed standard prerequisites (editorial).}\begin{itemize}
\item Filtered colimits of rings and modules, and their exactness. used throughout the counted proof span. Is what allows finite data over the colimit to be descended to a stage.
\item Perfect complexes as bounded complexes of finite projective modules. used in the counted proof span. Fixes the form in which the descending object is presented.
\end{itemize}\medskip
\ifdefined\hypertarget\hypertarget{sourcepointer-1}{}\fi
\noindent\textit{Source reference.} Original lemma 05N7: focused-seeds/AGP/sources/stacks-snapshot-20261011/tags.txt, source lines 7296--7296. Cross-chapter navigation pointer into the Stacks Project; the standard fact it supplies is applied inside the counted proof span and the argument resting on it stays there.\par\medskip
\ifdefined\hypertarget\hypertarget{sourcepointer-2}{}\fi
\noindent\textit{Source reference.} Original lemma 00DB: focused-seeds/AGP/sources/stacks-snapshot-20261011/tags.txt, source lines 491--491. Cross-chapter navigation pointer into the Stacks Project; the standard fact it supplies is applied inside the counted proof span and the argument resting on it stays there.\par\medskip
\ifdefined\hypertarget\hypertarget{sourcepointer-3}{}\fi
\noindent\textit{Source reference.} Original lemma 00NX: focused-seeds/AGP/sources/stacks-snapshot-20261011/tags.txt, source lines 875--875. Cross-chapter navigation pointer into the Stacks Project; the standard fact it supplies is applied inside the counted proof span and the argument resting on it stays there.\par\medskip
\ifdefined\hypertarget\hypertarget{sourcepointer-4}{}\fi
\noindent\textit{Source reference.} Original lemma 05LI: focused-seeds/AGP/sources/stacks-snapshot-20261011/tags.txt, source lines 7237--7237. Cross-chapter navigation pointer into the Stacks Project; the standard fact it supplies is applied inside the counted proof span and the argument resting on it stays there.\par\medskip
\ifdefined\hypertarget\hypertarget{sourcepointer-5}{}\fi
\noindent\textit{Source reference.} Original lemma 064B: focused-seeds/AGP/sources/stacks-snapshot-20261011/tags.txt, source lines 7860--7860. Cross-chapter navigation pointer into the Stacks Project; the standard fact it supplies is applied inside the counted proof span and the argument resting on it stays there.\par\medskip
\ifdefined\hypertarget\hypertarget{sourcepointer-6}{}\fi
\noindent\textit{Source reference.} Original lemma 14380: sources/stacks-more-algebra/pinned-original-source.tex, source lines 14380--14380. Same-chapter navigation pointer; the statement it names is standard for the section and the argument that consumes it stays inside the counted proof span.\par\medskip
\begin{lemma}
\label{lemma-colimit-perfect-complexes}
Let $R = \colim_{i \in I} R_i$ be a filtered colimit of rings.
\begin{enumerate}
\item Given a perfect $K$ in $D(R)$ there exists an $i \in I$
and a perfect $K_i$ in $D(R_i)$ such that
$K \cong K_i \otimes_{R_i}^\mathbf{L} R$ in $D(R)$.
\item Given $0 \in I$ and $K_0, L_0 \in D(R_0)$ with $K_0$ perfect,
we have
$$
\Hom_{D(R)}(K_0 \otimes_{R_0}^\mathbf{L} R, L_0 \otimes_{R_0}^\mathbf{L} R) =
\colim_{i \geq 0}
\Hom_{D(R_i)}(K_0 \otimes_{R_0}^\mathbf{L} R_i,
L_0 \otimes_{R_0}^\mathbf{L} R_i)
$$
\end{enumerate}
In other words, the triangulated category of perfect complexes over $R$
is the colimit of the triangulated categories of perfect complexes over $R_i$.
\end{lemma}

\begin{proof}
We will use the results of
Algebra, Lemmas \ref{algebra-lemma-module-map-property-in-colimit} and
\ref{algebra-lemma-colimit-category-fp-modules}
without further mention. These lemmas in particular say that the
category of finitely presented $R$-modules is the colimit of the
categories of finitely presented $R_i$-modules. Since finite projective
modules can be characterized as summands of finite free modules
(Algebra, Lemma \ref{algebra-lemma-finite-projective}) we see that
the same is true for the category of finite projective modules.
This proves (1) by our definition of perfect objects of $D(R)$.

\medskip\noindent
To prove (2) we may represent $K_0$ by a bounded complex $K_0^\bullet$ of
finite projective $R_0$-modules. We may represent $L_0$ by a K-flat
complex $L_0^\bullet$ (Lemma \ref{lemma-K-flat-resolution}).
Then we have
$$
\Hom_{D(R)}(K_0 \otimes_{R_0}^\mathbf{L} R, L_0 \otimes_{R_0}^\mathbf{L} R) =
\Hom_{K(R)}(K_0^\bullet \otimes_{R_0} R, L_0^\bullet \otimes_{R_0} R)
$$
by Derived Categories, Lemma
\ref{derived-lemma-morphisms-from-projective-complex}.
Similarly for the $\Hom$ with $R$ replaced by $R_i$. Since
in the right hand side only a finite number of terms are involved,
since
$$
\Hom_R(K_0^p \otimes_{R_0} R, L_0^q \otimes_{R_0} R) =
\colim_{i \geq 0}
\Hom_{R_i}(K_0^p \otimes_{R_0} R_i, L_0^q \otimes_{R_0} R_i)
$$
by the lemmas cited at the beginning of the proof, and since
filtered colimits are exact
(Algebra, Lemma \ref{algebra-lemma-directed-colimit-exact})
we conclude that (2) holds as well.
\end{proof}

\end{document}
